\section{Introduction}
\label{sec:introduction}
% Target: 1 page including title and abstract. Paragraph plan, not prose.
% Source mapping: content_mapping.md, Introduction.
\begin{itemize}
    \item \textbf{Paragraph 1---why structure matters.} Explain how task-specific
    computations and dependencies support learning from limited demonstrations.
    Introduce KIM briefly as prior work; keep its representation and empirical
    contributions distinct from this paper's refinement contribution.
    \item \textbf{Paragraph 2---why refinement is needed.} Preserve the motivation
    of reducing the expert's burden of specifying a strategy. Explain that an
    LLM's domain knowledge can differ from the strategy expressed in a
    demonstration, leaving poor behavior after parameter fitting. An initial
    structure is provisional, not necessarily incapable of solving the task.
    \item \textbf{Paragraph 3---semantics as an interface.} Connect meaningful
    intermediate quantities and explicit dependencies to measurable execution
    behavior and editable computations. Establish both roles of structure:
    inductive bias for fitting and an interface for refinement.
    \item \textbf{Paragraph 4---the iterative solution.} Preview the complete
    fit--execute--analyze--revise--refit cycle. Distinguish task descriptions
    used for initial generation from demonstrations used for parameter learning
    and policy rollouts used for diagnostic feedback.
    Use a brief debugging analogy: instrument execution, write diagnostic checks,
    inspect results, revise, and rerun. Explain that here the program is a learned
    controller, and each structural edit is followed by parameter fitting.
    \item \textbf{Paragraph 5---contributions and evidence.} State the iterative
    procedure, its semantic diagnostic interface, and the two-task refinement
    evaluation. Preview feedback comparisons and downstream RL as supporting
    evidence. Do not present the KIM representation as a new contribution.
\end{itemize}
